% =============================================================================
% ARBEITSBLATT: Fake News & KI-Bilder erkennen
% Informatik Klasse 10 MR
% =============================================================================

\documentclass[a4paper,11pt]{article}

\usepackage[utf8]{inputenc}
\usepackage[T1]{fontenc}
\usepackage[ngerman]{babel}
\usepackage{helvet}
\renewcommand{\familydefault}{\sfdefault}

\usepackage[margin=20mm]{geometry}
\usepackage{tikz}
\usepackage{xcolor}
\usepackage{tabularx}
\usepackage{enumitem}
\usepackage{amssymb}
\usepackage{qrcode}

% =============================================================================
% FARBEN
% =============================================================================

\definecolor{informatik}{HTML}{b35c00}
\definecolor{liniengrau}{HTML}{cccccc}
\definecolor{hintergrund}{HTML}{fff3e0}

% =============================================================================
% BEFEHLE
% =============================================================================

\newcommand{\schreiblinie}{\rule{\linewidth}{0.4pt}}
\newcommand{\kleinezeile}{\vspace{3mm}\hrulefill\vspace{3mm}}
\newcommand{\checkbox}{$\square$\;}

% =============================================================================
% DOKUMENT
% =============================================================================

\pagestyle{empty}

\begin{document}

% -----------------------------------------------------------------------------
% KOPFZEILE
% -----------------------------------------------------------------------------

\begin{tikzpicture}[remember picture, overlay]
    \fill[informatik] (current page.north west) rectangle ([yshift=-15mm]current page.north east);
    \node[anchor=west, white, font=\Large\bfseries] at ([xshift=20mm, yshift=-8mm]current page.north west) {Fake News \& KI-Bilder erkennen};
    \node[anchor=east, white, font=\small] at ([xshift=-20mm, yshift=-8mm]current page.north east) {Informatik 10MR};
\end{tikzpicture}

\vspace{8mm}

% -----------------------------------------------------------------------------
% NAMENSFELD MIT QR-CODE
% -----------------------------------------------------------------------------

\begin{tabularx}{\textwidth}{|X|l|c|}
\hline
\textbf{Name:} \hspace{4cm} & \textbf{Datum:} \hspace{2.5cm} & \raisebox{-4mm}{\qrcode[height=12mm]{https://devbox42.github.io/sciencesim/informatik/kl10-MR/01-fake-news-ki-bilder/LP-01-fake-news-ki-bilder.html}} \\[4mm]
\hline
\end{tabularx}

\vspace{5mm}

% -----------------------------------------------------------------------------
% AUFGABE 1: KI-BILDER ERKENNEN
% -----------------------------------------------------------------------------

\section*{Aufgabe 1: KI-Bilder erkennen}

Spiele 5 Runden auf \textbf{whichfaceisreal.com} und notiere deine Beobachtungen.

\vspace{3mm}

\textbf{Welche Merkmale haben dir geholfen, das KI-Bild zu erkennen?}

\vspace{2mm}

\begin{tabularx}{\textwidth}{|c|X|}
\hline
\textbf{Merkmal} & \textbf{Was hast du beobachtet?} \\
\hline
Ohren/Ohrringe & \vspace{8mm} \\
\hline
Hintergrund & \vspace{8mm} \\
\hline
Haare/Ränder & \vspace{8mm} \\
\hline
Sonstiges & \vspace{8mm} \\
\hline
\end{tabularx}

\vspace{5mm}

% -----------------------------------------------------------------------------
% AUFGABE 2: CRAAP-METHODE
% -----------------------------------------------------------------------------

\section*{Aufgabe 2: Quellen prüfen (CRAAP-Methode)}

Erkläre mit eigenen Worten, was die Buchstaben bedeuten:

\vspace{2mm}

\begin{tabularx}{\textwidth}{|c|l|X|}
\hline
& \textbf{Begriff} & \textbf{Deine Erklärung} \\
\hline
\textbf{C} & Currency (Aktualität) & \vspace{6mm} \\
\hline
\textbf{R} & Relevance (Relevanz) & \vspace{6mm} \\
\hline
\textbf{A} & Authority (Autorität) & \vspace{6mm} \\
\hline
\textbf{A} & Accuracy (Genauigkeit) & \vspace{6mm} \\
\hline
\textbf{P} & Purpose (Zweck) & \vspace{6mm} \\
\hline
\end{tabularx}

\vspace{5mm}

% -----------------------------------------------------------------------------
% AUFGABE 3: RECHERCHE
% -----------------------------------------------------------------------------

\section*{Aufgabe 3: Dein Fake-Beispiel}

Dokumentiere hier das Beispiel, das du für die Präsentation gefunden hast.

\vspace{3mm}

\textbf{Was hast du gefunden?} (Fake News / KI-Bild / Deepfake)

\vspace{2mm}
\fbox{\parbox{\dimexpr\linewidth-2\fboxsep-2\fboxrule}{\vspace{10mm}\hspace{1mm}}}

\vspace{4mm}

\textbf{Wo hast du es gefunden?} (URL oder Quelle)

\vspace{2mm}
\fbox{\parbox{\dimexpr\linewidth-2\fboxsep-2\fboxrule}{\vspace{8mm}\hspace{1mm}}}

\vspace{4mm}

\textbf{Woran hast du erkannt, dass es fake ist?}

\vspace{2mm}
\fbox{\parbox{\dimexpr\linewidth-2\fboxsep-2\fboxrule}{\vspace{15mm}\hspace{1mm}}}

\vspace{5mm}

% -----------------------------------------------------------------------------
% AUFGABE 4: REFLEXION
% -----------------------------------------------------------------------------

\section*{Aufgabe 4: Reflexion}

\textbf{Warum ist es gefährlich, wenn immer mehr Menschen Fake News nicht erkennen können? Welche Folgen kann das für die Gesellschaft haben?}

\vspace{2mm}
\fbox{\parbox{\dimexpr\linewidth-2\fboxsep-2\fboxrule}{\vspace{25mm}\hspace{1mm}}}

\vspace{8mm}

% -----------------------------------------------------------------------------
% SELBSTEINSCHÄTZUNG
% -----------------------------------------------------------------------------

\begin{tikzpicture}
\fill[hintergrund] (0,0) rectangle (\linewidth, 22mm);
\draw[informatik, thick] (0,0) rectangle (\linewidth, 22mm);
\node[anchor=west, font=\bfseries] at (3mm, 18mm) {Selbsteinschätzung:};
\node[anchor=west, font=\small] at (3mm, 12mm) {\checkbox Ich kann erklären, was Fake News sind.};
\node[anchor=west, font=\small] at (3mm, 6mm) {\checkbox Ich kann KI-generierte Bilder an typischen Fehlern erkennen.};
\node[anchor=west, font=\small] at (95mm, 12mm) {\checkbox Ich kann die CRAAP-Methode anwenden.};
\node[anchor=west, font=\small] at (95mm, 6mm) {\checkbox Ich habe ein Fake-Beispiel gefunden.};
\end{tikzpicture}

\end{document}
